\chapter{Huntington's Disease}
\index{Huntington's Disease}

\section*{Definition and Overview}
Huntington's disease is a rare, inherited neurodegenerative disorder that causes progressive damage to the brain, leading to movement problems, cognitive decline, and psychiatric symptoms. It is caused by an abnormal gene on chromosome 4, which produces a toxic form of the huntingtin protein. Symptoms usually begin in mid-adulthood, typically between 30 and 50 years of age, and worsen gradually over 10--25 years. Each child of an affected parent has a 50\% chance of inheriting the gene. There is no cure, but treatments manage symptoms and improve quality of life, and research into disease-modifying therapies is active.

\section*{Epidemiology}
Huntington's disease affects around 1 in 7,000--10,000 people worldwide, with roughly 2,000--3,000 people affected and a similar number at risk. Onset is most common between 30 and 50, but juvenile Huntington's disease occurs in about 5--10\% of cases, before age 20. The condition affects men and women equally. Predictive genetic testing allows people at risk to learn whether they carry the gene, and around 10--20\% of at-risk adults choose to be tested.

\section*{Aetiology and Pathophysiology}
\begin{itemize}
    \item \textbf{Genetic cause:} an expanded CAG repeat in the HTT gene on chromosome 4
    \item \textbf{Inheritance:} autosomal dominant; each child of an affected parent has a 50\% risk
    \item \textbf{Repeat length:} more than 40 CAG repeats causes the disease; longer repeats cause earlier onset
    \item \textbf{Pathophysiology:} the expanded gene produces a mutant huntingtin protein that is toxic to nerve cells, particularly in the basal ganglia and cortex
    \item \textbf{Cell loss:} progressive death of neurons causes the movement, cognitive, and psychiatric features
    \item \textbf{Juvenile onset:} large repeat expansions cause childhood disease with different features
    \item \textbf{Anticipation:} repeat length can increase across generations, causing earlier onset in some families
\end{itemize}

\section*{Clinical Features}
\begin{itemize}
    \item Movement problems: involuntary jerking (chorea), clumsiness, and difficulty with coordination
    \item Rigidity, slow movements, and balance problems in later stages
    \item Cognitive decline: impaired memory, planning, and judgement
    \item Psychiatric symptoms: depression, anxiety, irritability, and obsessive behaviour
    \item Personality changes and social withdrawal
    \item Difficulty with speech and swallowing in later stages
    \item Weight loss and sleep disturbance
    \item In juvenile forms: rigidity, seizures, and learning decline
\end{itemize}

\section*{Differential Diagnosis}
\begin{itemize}
    \item Other causes of chorea, including drug-induced and metabolic causes
    \item Other neurodegenerative conditions, including Parkinson's disease and dementias
    \item Depression and other psychiatric disorders, which may precede movement symptoms
    \item Genetic testing confirms the diagnosis
    \item Other repeat-expansion disorders are considered in atypical cases
\end{itemize}

\section*{When to Seek Medical Care}
Anyone with progressive movement, cognitive, or psychiatric symptoms should be assessed by a neurologist, and people with a family history can seek genetic counselling and predictive testing.

\textbf{Contact your local emergency services for:}
\begin{itemize}
    \item Suicidal thoughts or self-harm
    \item Choking or severe difficulty swallowing
    \item Falls with suspected head injury or fracture
    \item Severe agitation or confusion
    \item Any life-threatening emergency
\end{itemize}

\section*{Diagnosis and Investigations}
\begin{itemize}
    \item \textbf{Clinical assessment:} history and neurological examination
    \item \textbf{Genetic testing:} confirms the diagnosis by measuring the CAG repeat length
    \item \textbf{Predictive testing:} available for at-risk adults through specialist programs with counselling
    \item \textbf{Pre-symptomatic testing:} involves careful psychological preparation and support
    \item \textbf{MRI:} may show atrophy of the basal ganglia
    \item \textbf{Prenatal testing:} options exist for families planning children
\end{itemize}

\section*{Management}
\begin{itemize}
    \item \textbf{Symptom management:} medicines for chorea, movement, and psychiatric symptoms
    \item \textbf{Psychiatric care:} treatment of depression, anxiety, and irritability
    \item \textbf{Physiotherapy:} maintain mobility, balance, and function
    \item \textbf{Speech and swallowing therapy:} manage communication and reduce aspiration risk
    \item \textbf{Occupational therapy:} adapt daily activities and maintain independence
    \item \textbf{Multidisciplinary care:} neurologists, psychiatrists, nurses, and allied health
    \item \textbf{Support services:} counselling, support groups, and community care
    \item \textbf{Research:} clinical trials of disease-modifying therapies
\end{itemize}

\section*{Complications}
\begin{itemize}
    \item Falls, fractures, and head injuries
    \item Choking and aspiration pneumonia from swallowing difficulty
    \item Malnutrition and weight loss
    \item Depression and suicide risk
    \item Progressive dependence on care
    \item Pneumonia and infections in advanced disease
    \item Psychological impact on families and caregivers
\end{itemize}

\section*{Prognosis and Outcomes}
Huntington's disease is progressive, with symptoms worsening over 10--25 years from onset, and there is currently no cure. However, multidisciplinary care improves quality of life, and symptom treatments are effective for many features. Disease-modifying therapies are in clinical trials and offer hope for the future. Families benefit from genetic counselling, predictive testing support, and comprehensive care services.

\section*{Prevention and Self-Care}
\begin{itemize}
    \item Seek genetic counselling if you have a family history
    \item Consider predictive testing only with full support and preparation
    \item Discuss reproductive options with a genetics service
    \item Maintain physical activity and nutrition as long as possible
    \item Attend specialist clinics and allied health services
    \item Join support groups for affected people and families
    \item Plan for future care needs early
    \item Caregivers should seek respite and support
\end{itemize}

\section*{Sources and Further Reading}
The following reputable sources provide further information for healthcare students and patients seeking reliable, current advice about Huntington's disease.

\begin{itemize}
    \item Huntington's Australia. \textit{Information and support.} https://huntingtonsaustralia.org.au/
    \item Healthdirect. \textit{Huntington's disease.} https://www.healthdirect.gov.au/huntingtons-disease
    \item Brain Foundation. \textit{Huntington's disease information.} https://brainfoundation.org.au/
    \item NHS. \textit{Huntington's disease: overview.} https://www.nhs.uk/conditions/huntingtons-disease/
    \item Mayo Clinic. \textit{Huntington's disease: symptoms and causes.} https://www.mayoclinic.org/
    \item Huntington's Disease Society of America. \textit{Information and research.} https://hdsa.org/
\end{itemize}
